% Part: normal-modal-logic
% Chapter: frame-definability
% Section: properties-accessibility

\documentclass[../../../include/open-logic-section]{subfiles}

\begin{document}

\olfileid{nml}{frd}{acc}

\olsection{સુલભતા સંબંધોના ગુણધર્મો}

ઘણાં મોડલ !!{formula}s સુલભતા સંબંધના સરળ, અને
પરિચિત પણ, ગુણધર્મોને લક્ષણરૂપ નીવડે છે. એક
દિશામાં તેનો અર્થ એ છે કે જે નિદર્શનનો સંબંધ
આપેલો ગુણધર્મ ધરાવે છે તેમાં અનુરૂપ !!{formula}
(અને તેના બધા પ્રતિસ્થાપન દૃષ્ટાંતો) સત્ય હોય છે.
આપણે સુલભતા સંબંધોના પાંચ શાસ્ત્રીય પ્રકારો
અને તેઓ જે સૂત્રોની સત્યતા સુનિશ્ચિત કરે છે
તેનાથી શરૂઆત કરીએ.

\begin{thm}\ollabel{thm:soundschemas}
  $\mModel{M} = \tuple{W, R, V}$ નિદર્શન લો.
  જો $R$ એ \olref{tab:five}ની ડાબી બાજુનો ગુણધર્મ
  ધરાવે, તો જમણી બાજુના !!{formula}નો દરેક
  દૃષ્ટાંત~$\mModel{M}$માં સત્ય છે.
\end{thm}

\begin{table}[t]
    \begin{tabular}{| l || l |}
      \hline
      {\emph{જો $R$ \dots હોય}} & {\emph{તો \dots એ~$\mModel{M}$માં સત્ય છે:}} \\
      \hline \hline
      \emph{સિરિયલ}: $\forall u \exists v Ruv$ & \hbox to.43\textwidth
           {$\Box p \lif \Diamond p$ \hfill (\Ax{D})} \\
      \hline
      \emph{સ્વવાચક}: $\forall w Rww$  
      & \hbox to.43\textwidth
          {$\Box p \lif p$ \hfill (\Ax{T})} \\
      \hline
      \emph{સંમિત}: &  \hbox to .43\textwidth
           {$p \lif \Box\Diamond p$ \hfill (\Ax{B})} \\
           $\forall u\forall v(Ruv \lif Rvu)$ & \\
      \hline
      \emph{પરંપરિત}: & \hbox to.43\textwidth
           {$\Box p \lif \Box \Box p$ \hfill (\Ax{4})} \\
      $\forall u \forall v \forall w ((Ruv \land Rvw) \lif Ruw)$ & \\
      \hline 
      \emph{યુક્લિડિયન}: & \hbox to .43\textwidth
        {$\Diamond p \lif \Box\Diamond p$ \hfill (\Ax{5})} \\
      $\forall w \forall u \forall v ((Rwu \land Rwv) \lif Ruv)$ &\\
      \hline
    \end{tabular}
    \caption{અનુરૂપતાનાં પાંચ તથ્યો.}
    \ollabel{tab:five}
  \end{table} 


\begin{proof}
  અહીં \Ax{B}નો કિસ્સો જોઈએ: સૂત્ર-ઢાંચો નિદર્શનમાં
  સત્ય છે તે બતાવવા તેના બધા દૃષ્ટાંતો નિદર્શનના બધા
  જગતોમાં સત્ય છે તે બતાવવું પડે. તેથી \Ax{B}નો
  કોઈ દૃષ્ટાંત $!A \lif \Box\Diamond !A$ લો અને
  $w \in W$ કોઈપણ જગત લો. $\Box \Diamond
  !A$ એ $w$માં સત્ય છે તે બતાવવા $!A$ પૂર્વપક્ષ
  $w$માં સત્ય છે એમ માનો. એટલે $w$માંથી સુલભ
  દરેક $w'$માં $\Diamond !A$ સત્ય છે તે બતાવવું છે.
  હવે $Rww'$ ધરાવતા કોઈપણ $w'$ માટે સંમિતતાની
  ધારણાથી $Rw'w$ પણ છે (જુઓ \olref{fig:Bsymm}).
  $\mSat{M}{!A}[w]$ હોવાથી
  $\mSat{M}{\Diamond !A}[w']$ છે. $Rww'$ ધરાવતું
  $w'$ ગમે તે હતું, તેથી
  $\mSat{M}{\Box\Diamond!A}[w]$.

  બાકીના કિસ્સા અભ્યાસપ્રશ્ન તરીકે છોડીએ છીએ.
\end{proof}

\begin{prob}
  \olref[nml][frd][acc]{thm:soundschemas}નું પ્રમાણ
  પૂરું કરો.
\end{prob}

\begin{figure}
  \begin{center}
    \begin{tikzpicture}[modal]
      \node[world] (w1) [label={below:$\mSat{{}}{\formula{A}}$\\
          $\mSat{{}}{\Box\Diamond \formula{A}}$}] {$w$}; 
      \node[world] (w2) [label=below:{$\mSat{{}}{\Diamond \formula{A}}$},
        right=of w1] {$w'$}; 
      \draw[->, bend left] (w1) to (w2); 
      \draw[->, bend left] (w2) to (w1); 
    \end{tikzpicture}
  \end{center}
\caption{સંમિતતાથી થતી દલીલ.}
\ollabel{fig:Bsymm}
\end{figure}

ધ્યાન આપો કે \olref{thm:soundschemas}નાં વિપરીત
વિધાનો સાચાં નથી: કોઈ નિદર્શનમાં સૂત્ર-ઢાંચો
સત્ય હોય તો તે નિદર્શનનો સુલભતા સંબંધ અનુરૂપ
ગુણધર્મ ધરાવે જ એવું નથી. \Ax{T} અને સ્વવાચક
નિદર્શનોના કિસ્સામાં \Ax{T} પોતે નિદર્શનમાં ખોટું
પડે એવો દાખલો સરળતાથી મળે: $W = \{w\}$ અને
$V(p) = \emptyset$ લો, તથા સુલભતા સંબંધ ખાલી લો.
ત્યારે $R$ સ્વવાચક નથી, પણ
$\mSat{M}{\Box
  p}[w]$ અને $\mSat/{M}{p}[w]$. જોકે અહીં
\Ax{T}નો માત્ર એક દૃષ્ટાંત~$\mModel{M}$માં
ખોટો છે; $\Box \lnot p \lif \lnot p$ જેવા
બીજા દૃષ્ટાંતો સાચા છે. \Ax{T}નો
\emph{દરેક પ્રતિસ્થાપન દૃષ્ટાંત}~$\mModel{M}$માં
સત્ય હોય અને $\mModel{M}$ સ્વવાચક ન હોય એવાં
નિદર્શનો આપવાં વધુ મુશ્કેલ છે. પરંતુ એવાં
નિદર્શનો પણ છે:
\sourcecorrection{OLMD-013}{સ્થિર અંગ્રેજી મૂળના
એક-જગત વિરોધી નિદર્શનમાં સુલભતા સંબંધનો
ઉલ્લેખ કર્યા વિના તે અસ્વવાચક હોવાનું કહે છે.
અહીં ખાલી સંબંધ સ્પષ્ટ કર્યો છે; નહિતર દાવો
આપેલી માહિતીમાંથી ફલિત થતો નથી.}

\begin{prop}\ollabel{prop:reflexive}
  $\mModel{M} = \tuple{W, R, V}$ એવું નિદર્શન
  લો કે $W = \{u, v
  \}$ અને જગતો $u$ તથા $v$ સંબંધ $R$થી
  જોડાયેલાં છે: એટલે કે $Ruv$ અને $Rvu$ બંને છે;
  આ બે સિવાય સંબંધમાં બીજી કોઈ જોડી નથી.
  વળી, બધા $p$ માટે $u \in V(p) \Leftrightarrow v
  \in V(p)$ માનો. ત્યારે:
  \begin{enumerate}
  \item બધા $!A$ માટે: $\mSat{M}{!A}[u]$ ત્યારે અને માત્ર ત્યારે
    જ જ્યારે $\mSat{M}{!A}[v]$ (અહીં $!A$ પર આગમન વાપરો).
  \item \Ax{T}નો દરેક દૃષ્ટાંત $\mModel{M}$માં સત્ય છે.
  \end{enumerate}
  $\mModel{M}$ સ્વવાચક નથી (હકીકતમાં તે
  \emph{અસ્વવાચક} છે), તેથી \Ax{T}ના કિસ્સામાં
  \olref{thm:soundschemas}નું વિપરીત વિધાન ખોટું છે
  (\olref{thm:soundschemas}માં જણાવેલા અન્ય કેટલાંક---જોકે બધાં નહીં---
  સૂત્ર-ઢાંચાઓ માટે પણ આવી જ દલીલો આપી શકાય).
\end{prop}
\sourcecorrection{OLMD-014}{સ્થિર અંગ્રેજી મૂળમાં
માત્ર બંને આડાં તીરોની ધારણા છે, છતાં અંતે
નિદર્શનને અસ્વવાચક કહે છે. સ્વ-તીરોને બહાર રાખતી
શરત સ્પષ્ટ કરી છે; તે વિના અંતિમ નિષ્કર્ષ આવતો નથી.}

\begin{prob}
  \olref[nml][frd][acc]{prop:reflexive}ના દાવા સાબિત કરો.
\end{prob}

આપણે પાંચ શાસ્ત્રીય !!{formula}s \Ax{D},
\Ax{T}, \Ax{B}, \Ax{4} અને \Ax{5} પર ધ્યાન કેન્દ્રિત
કરીશું; છતાં \olref{tab:anotherfive}માં સુલભતા
સંબંધોના થોડા વધુ ગુણધર્મો નોંધીએ. સુલભતા
સંબંધ~$R$ \emph{આંશિક વિધેયાત્મક} છે જો દરેક
જગતમાંથી વધુમાં વધુ એક જગત સુલભ હોય. દરેક
જગતમાંથી બરાબર એક જગત સુલભ હોય તો તેને
\emph{વિધેયાત્મક} કહીએ. (આમ, વિધેયાત્મક સંબંધો
બરાબર સિરિયલ અને આંશિક વિધેયાત્મક બંને હોય
તેવા સંબંધો છે.) તેમને ``વિધેયાત્મક'' કહે છે કારણ કે
સુલભતા સંબંધ (આંશિક) વિધેયની જેમ વર્તે છે.
સંબંધ \emph{નબળે સઘન} છે જો $Ruv$ હોય ત્યારે
$u$ અને~$v$ની ``વચ્ચે'' કોઈ~$w$ હોય. તેથી
નબળે સઘન સંબંધો એક અર્થમાં પરંપરિત સંબંધોથી
વિપરીત છે: પરંપરિત સંબંધમાં~$u$માંથી~$w$ થઈને
$v$ સુધી પહોંચી શકાય તો~$u$માંથી સીધા~$v$
સુધી પણ પહોંચી શકાય; નબળે સઘન સંબંધમાં
~$u$માંથી સીધા~$v$ સુધી પહોંચી શકાય તો
કોઈ~$w$ થઈને પણ ત્યાં પહોંચી શકાય.
સંબંધ \emph{નબળે દિશિત} છે જો કોઈ જગત~$w$માંથી
$u$ અને~$v$ બંને સુધી પહોંચી શકાય ત્યારે~$u$
અને~$v$ બંનેમાંથી એક જ જગત~$t$ સુધી પહોંચી
શકાય---આને ક્યારેક ``હીરા ગુણધર્મ'' અથવા
``અભિસરણ'' કહે છે.

\begin{table}[t]
    \begin{tabular}{| p{.50\textwidth} || p{.43\textwidth} |}
      \hline
      {\emph{જો $R$ \dots હોય}} & {\emph{તો \dots એ~$\mModel{M}$માં સત્ય છે:}} \\
      \hline \hline
      \emph{આંશિક વિધેયાત્મક}: & 
      \multirow{2}{*}{$\Diamond p \lif \Box p$} \\
      $ \forall w \forall u \forall v ((Rwu \land Rwv) \lif u =v )$ & \\
      \hline
      \multirow{2}{*}{\emph{વિધેયાત્મક}: $\forall w \exists v \forall u(Rwu \liff \eq[u][v]) $} 
      & \multirow{2}{*}{$\Diamond p \liff \Box p$} \\
      & \\
      \hline
      \emph{નબળે સઘન}: &  \multirow{2}{*}{$\Box \Box p \lif
        \Box p$} \\
      $\forall u\forall v(Ruv \lif \exists w(Ruw \land Rwv))$ & \\
      \hline
      \emph{નબળે જોડાયેલો}: &
      \multirow{3}{*}{\hbox to.43\textwidth{$
        \begin{array}{@{}l@{}}
          \Box ((p \land \Box p) \lif q) \lor {} \\
          \qquad \Box ((q \land \Box q) \lif p)
        \end{array}$ \hfill (\Ax{L})}
      } \\
      $\forall w \forall u \forall v ((Rwu \land Rwv) \lif {}$ &\\
      \qquad $(Ruv \lor u=v \lor Rvu))$ & \\
      \hline 
      \emph{નબળે દિશિત}: & \multirow{3}{*}{\hbox to .43\textwidth
        {$\Diamond\Box p \lif \Box\Diamond p$\hfill (\Ax{G})}} \\
      $\forall w \forall u \forall v ((Rwu \land Rwv) \lif {}$ & \\
      \qquad $\exists t (Rut \land Rvt))$ & \\
      \hline
    \end{tabular}
    \caption{અનુરૂપતાનાં વધુ પાંચ તથ્યો.}
    \ollabel{tab:anotherfive}
  \end{table} 

\begin{prob}
  $\mModel{M} = \tuple{W, R, V}$ નિદર્શન લો.
  જો $R$ એ \olref[nml][frd][acc]{tab:anotherfive}ની
  ડાબી બાજુના ગુણધર્મો ધરાવતું હોય, તો અનુરૂપ
  જમણી બાજુના સૂત્રનો દરેક દૃષ્ટાંત~$\mModel{M}$માં
  સત્ય છે તે બતાવો.
\end{prob}

\end{document}
